% =====================================================================
% บทที่ 2 งานวิจัยและทฤษฎีที่เกี่ยวข้อง (คู่มือข้อ 17.2)
%   2.1 งานวิจัย โครงงาน ระบบงาน หรือโปรแกรมที่เกี่ยวข้อง
%       (ส่วนที่คล้ายกัน และส่วนที่แตกต่างจากโครงงานนี้)
%   2.2 หลักการและทฤษฎีที่นำมาใช้ในโครงงาน
%
% การอ้างอิง: \cite{key} -> [1]   \cite{a,b} -> [1, 2]   \cite{a,b,c} -> [1-3]
% เลขอ้างอิงเรียงตามลำดับที่อ้างครั้งแรกในเล่ม และแสดงในหน้า "เอกสารอ้างอิง" ให้อัตโนมัติ
% *** ก่อนส่ง: ลบข้อความตัวอย่าง (...) ***
% =====================================================================

\section{งานวิจัย โครงงาน และระบบงานที่เกี่ยวข้อง}
\label{sec:related}

% สำหรับแต่ละงาน: ใครทำ ทำอะไร ใช้วิธีใด ได้ผลอย่างไร จุดเด่น/ข้อจำกัด
% และเกี่ยวข้องกับโครงงานนี้อย่างไร (นำแนวคิดใดมาใช้ หรือโครงงานนี้แก้ข้อจำกัดใด)

ผู้จัดทำได้ศึกษางานวิจัยและระบบที่มีลักษณะคล้ายกับโครงงานนี้ จำนวน 4 งาน
ได้แก่ ระบบจัดการวิทยานิพนธ์ออนไลน์~\cite{thesismgmt2023}
ระบบจัดการโครงงานสำหรับรายวิชาวิศวกรรมซอฟต์แวร์~\cite{projmgmt2023}
และโปรแกรมแก้ไขเอกสาร \LaTeX{} ออนไลน์~\cite{overleaf2024,swiftlatex2019}
มีรายละเอียดดังนี้

\subsection{ระบบจัดการวิทยานิพนธ์ออนไลน์}

สมหญิง รักเรียน และสมศักดิ์ ขยันทำ~\cite{thesismgmt2023} พัฒนาระบบ\ldots{}
โดยใช้\ldots{} ผลการทดสอบพบว่า\ldots{}
จุดเด่นของระบบคือ\ldots{} แต่ยังมีข้อจำกัดด้าน\ldots{}
โครงงานนี้นำแนวคิด\ldots{} มาใช้ และเพิ่มความสามารถ\ldots{}

\subsection{ระบบจัดการโครงงานสำหรับรายวิชาวิศวกรรมซอฟต์แวร์}

Somchai และ Wichian~\cite{projmgmt2023} นำเสนอ\ldots{}
ซึ่งคล้ายกับโครงงานนี้ในส่วน\ldots{} แต่แตกต่างกันที่\ldots{}

\subsection{สรุปการเปรียบเทียบ}

การเปรียบเทียบความสามารถของงานที่เกี่ยวข้องกับโครงงานนี้แสดงดังตารางที่~\ref{tab:compare}

\begin{table}[H]
\caption{การเปรียบเทียบความสามารถของงานที่เกี่ยวข้องกับโครงงานนี้}
\label{tab:compare}
\centering
\begin{tabular}{|L{4.6cm}|C{1.9cm}|C{1.9cm}|C{1.9cm}|C{1.9cm}|}
\hline
\textbf{งาน} & \textbf{ความสามารถ A} & \textbf{ความสามารถ B}
  & \textbf{ความสามารถ C} & \textbf{ความสามารถ D} \\ \hline
ระบบจัดการวิทยานิพนธ์~\cite{thesismgmt2023} & \cmark & \cmark & -- & -- \\ \hline
ระบบจัดการโครงงาน~\cite{projmgmt2023}      & \cmark & --     & \cmark & -- \\ \hline
Overleaf~\cite{overleaf2024}                 & --     & \cmark & \cmark & -- \\ \hline
\textbf{โครงงานนี้}                          & \cmark & \cmark & \cmark & \cmark \\ \hline
\end{tabular}
\end{table}

จากตารางที่~\ref{tab:compare} ยังไม่มีงานใดที่มีความสามารถ D
โครงงานนี้จึงพัฒนาระบบที่มีความสามารถครบทั้ง 4 ด้าน

\section{หลักการและทฤษฎีที่เกี่ยวข้อง}
\label{sec:theory}

% อธิบายเฉพาะหลักการ ทฤษฎี และเทคโนโลยีที่ "ใช้จริง" ในโครงงาน พร้อมอ้างอิง
% ไม่ต้องอธิบายความรู้ทั่วไปที่ไม่ได้นำมาใช้

\subsection{สถาปัตยกรรมแบบ REST}

REST (Representational State Transfer) เป็นรูปแบบสถาปัตยกรรมของระบบบนเครือข่าย
ที่นำเสนอโดย Fielding~\cite{fielding2000} มีหลักการสำคัญ ได้แก่
\begin{enumerate}
  \item การแยกส่วนผู้ใช้บริการและผู้ให้บริการ (Client--Server)
  \item การไม่เก็บสถานะของผู้ใช้บริการไว้ที่ผู้ให้บริการ (Stateless)
  \item การใช้ส่วนต่อประสานที่เป็นรูปแบบเดียวกัน (Uniform interface)
\end{enumerate}
โครงงานนี้ใช้หลักการดังกล่าวออกแบบ API ของระบบในหัวข้อ~\ref{sec:design}

\subsection{ระบบจัดเตรียมเอกสาร \LaTeX}

\LaTeX{} เป็นระบบจัดเตรียมเอกสารที่แยกเนื้อหาออกจากรูปแบบ~\cite{latex2023}
\ldots{}

\subsection{คอนเทนเนอร์ด้วย Docker}

Docker เป็นแพลตฟอร์มสำหรับสร้างและรันแอปพลิเคชันในคอนเทนเนอร์~\cite{docker2024}
\ldots{}
